\begin{abstract}
How much of the disagreement among theories of consciousness reflects genuine empirical differences, and how much is merely ``aesthetic''---redescribing the same behavioral facts in different representational vocabularies?
We propose the \emph{no costume theories} rule: any theoretical distinction that does not yield a behavioral test is empirically idle.
To operationalize this principle, we introduce three quantitative metrics---\costumeratio, prediction agreement, and \costumesens---and apply them to five major theories of consciousness (Integrated Information Theory, Global Neuronal Workspace, Higher-Order Thought, Recurrent Processing Theory, and Attention Schema Theory) using \gptmodel as a calibrated proxy reasoner across three experiments (${\sim}175$ evaluations).
We find that \IIT is overwhelmingly a ``costume theory'': 92\% of its predictions are representational-only with no behavioral test, compared to 20--30\% for the other four theories ($p = 0.007$, Kruskal--Wallis).
When behavior is held constant but substrate changes, \IIT consciousness attributions drop by 84 points on a 0--100 scale (Cohen's $d = 5.09$), nearly double the drop of any other theory.
Across 20 ambiguous scenarios, \IIT agrees with other theories only 35--45\% of the time, while the remaining four theories form a cohesive cluster (65--87\% agreement).
These results provide the first systematic quantification of ``costume-ness'' in consciousness science, identifying \IIT as the theory most vulnerable to the behavioral-testability criterion and offering transferable metrics for evaluating empirical content in any theoretical framework.
\end{abstract}
